\documentclass[10pt,a4paper]{amsart}
\usepackage[margin=19mm]{geometry}
\usepackage{amsmath,amssymb,mathtools}
\usepackage[scr=rsfs,cal=euler]{mathalfa}
\usepackage{xfrac}
\usepackage[unicode,hidelinks,bookmarks=false]{hyperref}
\newcommand{\ie}{i.e.\ }
\newcommand{\cf}{cf.\ }
\newcommand{\resp}{respectively\ }
\newcommand{\Subsection}{\subsection}
\providecommand{\nobreakdash}{\nobreak\hbox{-}}
\setlength{\emergencystretch}{2em}
\multlinegap0pt
\usepackage{bm}
\usepackage{bbm}
\usepackage{dsfont}
\usepackage{xspace}
\usepackage{cancel}
\usepackage{mathtools}
\usepackage{amsthm}
\makeatletter
\@ifundefined{theorem}{\newtheorem{theorem}{Theorem}}{}
\@ifundefined{proposition}{\newtheorem{proposition}[theorem]{Proposition}}{}
\makeatother
\title[Composition and tensor train structure in polynomial optimiz]{Composition and tensor train structure in polynomial optimization}
\author{Lloren\c{c} Balada Gaggioli, Didier Henrion, Milan Korda}
\date{Source: arXiv:2604.17563v1 (CC BY 4.0)}
\begin{document}
\maketitle

\noindent\emph{Source and licence.} Composition and tensor train structure in polynomial optimization. Version arXiv:2604.17563v1 (CC BY 4.0), distributed under the Creative Commons Attribution 4.0 International licence, as recorded in the arXiv licence metadata for this exact version and shown on the abstract page https://arxiv.org/abs/2604.17563v1. Full rights notice: licenses/arxiv-2604.17563-CC-BY-4.0.txt.\par\smallskip

\noindent\textit{Editorial scope.} Counted unit: the Proposition whose source label is \texttt{prop:gluing} in the article (source0.tex lines 1121--1124, source label \texttt{prop:gluing}), delivered together with the article's own complete original proof of it (source0.tex lines 1126--1150, one \texttt{proof} environment of 25 lines). No mathematics is written, repaired or re-derived: statement and proof are the article's own text line for line. Required context retained uncounted: none. Every definition, display and label of those lines is retained unchanged. Omitted from this package and not redistributed: 1--1120, 1125--1125, 1151--1850. Nothing inside a retained excerpt is omitted or altered: every retained body is delivered exactly as the article's lines are written, so no omission receipt of any kind accompanies this package. Citations: the retained excerpts carry no citation command at all, so no bibliography entry of the article is transcribed and no reference is redistributed. Reference adaptation: none was needed and none was made; every reference inside the delivered unit resolves inside it, so no unresolved reference or citation is produced. Printed slips kept literal: no printed word, symbol, hypothesis or number of the counted statement or of its proof was altered. Layout and class: the article's own class is \texttt{article} with packages including geometry, inputenc, subfiles, mathtools, esdiff, array, gensymb, amssymb, amsmath, bm, amsthm, amsfonts, braket, nth; the wrapper is the lane's pinned \texttt{amsart} editorial wrapper at 10pt on A4, which restates the theorem-like environments the retained text needs and adds the article's own macro definitions that the retained text needs (none), with extra wrapper packages (bm, bbm, dsfont, xspace, cancel, mathtools). The delivered numbering is the unit's own. Assets: no retained span contains a figure environment, a graphics-inclusion command, a PDF-page inclusion, a table or a TikZ picture; no image file is copied into this package, the package declares no asset dependency and no image file is redistributed with it. Licence: version arXiv:2604.17563v1 of this article is distributed under the Creative Commons Attribution 4.0 International licence, as recorded in the arXiv licence metadata for this exact version and shown on the abstract page https://arxiv.org/abs/2604.17563v1. A full rights notice is delivered at licenses/arxiv-2604.17563-CC-BY-4.0.txt. Characters: every retained excerpt is pure ASCII, so the delivered engine, XeLaTeX with its default Latin Modern text font, reports no missing character.
\begin{proposition}
Assume each local feasible set $K_i$ is compact. Then the infinite-dimensional $\mathrm{SL}_{\mathrm{push}}$ measure formulation has the same optimal value as the original compositional POP.
\label{prop:gluing}
\end{proposition}

\begin{proof}
We prove both inequalities.

First, let $(x_1,\dots,x_n)$ be any feasible point of the original problem, and define the states at each step recursively by
\[
s_1=F_1(x_1), \qquad s_i=F_i(s_{i-1},x_i), \quad i=2,\dots,n.
\]
This feasible trajectory induces a family of Dirac measures:
\[
\mu_1=\delta_{x_1}, \quad \mu_i=\delta_{(s_{i-1},x_i)}.
\]
for $i=2,\dots, n$. These measures satisfy the $\mathrm{SL}_{\mathrm{push}}$ support, marginal consistency constraints, and the objective value is exactly $s_n$. Therefore, the $\mathrm{SL}_{\mathrm{push}}$ measure optimum is at most the optimum of the original problem.

Conversely, let $(\mu_i)_{i=1}^n$ be any feasible family for the $\mathrm{SL}_{\mathrm{push}}$ measure formulation. By the marginal consistency constraints,
\[
(\pi_{s_{i-1}})_\#\mu_i\;=\; (F_{i-1})_{\#}\mu_{i-1}, \qquad i=2,\dots,n.
\]
Thus the local measures are consistent on their overlaps. The problem has a chain structure, and the overlap of all local blocks is exactly the state variable $s_{i-1}$, therefore by imposing the marginal consistency constraints we are ensuring that these shared variables are the same if we look at them from one stage or from the next one. Specifically, we have $\mu_1(x_1)$ and $\mu_2(s_1,x_2)$ with the marginals satisfying $(\pi_{s_{1}})_\#\mu_2= (F_{1}(x_1))_{\#}\mu_{1}$. Now we can define the image of $\mu_1$ via the map $x_1\mapsto (x_1,F_1(x_1))$ to define the auxiliary measure $\hat{\mu}_1(x_1,s_1)$. Similarly, via the map $(s_1,x_2)\mapsto (s_1,x_2,F_2(s_1,x_2))$ we define the auxiliary measure $\hat{\mu}_2(s_1,x_2,s_2)$, where $(\pi_{s_{1}})_\#\hat{\mu}_1= (\pi_{s_{1}})_\#\hat{\mu}_2$, which allows to glue them through their overlap $s_1$, to get the measure $\hat{\mu}_{(1,2)}(x_1,s_1,x_2,s_2)$. 


We can repeat this process consecutively, the next step would be taking $\hat{\mu}_{(1,2)}(x_1,s_1,x_2,s_2)$ and $\hat{\mu}_3(s_2,x_3,s_3)$, satisfying the marginal condition $(\pi_{s_{2}})_\#\hat{\mu}_{(1,2)}= (\pi_{s_{2}})_\#\hat{\mu}_3$, which yields a new measure $\hat{\mu}_{(1,2,3)}(x_1,s_1,x_2,s_2,x_3,s_3)$. After repeating this process a total of $n-1$ times we obtain a global measure $\hat{\mu}_T(x_1,s_1,x_2,\dots,s_n)$.


Moreover, by the support and push-forward constraints, this global measure is supported on the feasible set of the original compositional problem. Hence, the $\mathrm{SL}_{\mathrm{push}}$ objective is an average of feasible objective values of the original problem, so it is bounded below by the global optimum. Therefore, the $\mathrm{SL}_{\mathrm{push}}$ measure optimum is at least the optimum of the original problem. Combining both inequalities proves the claim.
\end{proof}

\end{document}
